\SourceLeaf{647}
\section*{INHALT.}
\addcontentsline{toc}{section}{Inhalt}
\begin{center}\bfseries GAUSS WERKE BAND V.\par MATHEMATISCHE PHYSIK.\end{center}
\par\nopagebreak[4]\smallskip\noindent\makebox[\linewidth][c]{\rule{.075\linewidth}{.35pt}}\par
\begingroup\fontsize{9}{11}\selectfont\raggedright
\par\Needspace{5\baselineskip}\medskip\noindent\textit{Abhandlungen}\par\smallskip
\noindent\hangindent=1em Theorie der Anziehung homogener ellipsoidischer Körper \unskip\hspace{1em plus 1fill}\mbox{1813 März . .\quad Seite 1}\par\smallskip
\noindent\hangindent=1em Über ein neues allgemeines Grundgesetz der Mechanik \unskip\hspace{1em plus 1fill}\mbox{1829\quad — 23}\par\smallskip
\noindent\hangindent=1em Allgemeine Grundlagen der Theorie der Gestalt von Flüssigkeiten im Gleichgewicht . . .\unskip\hspace{1em plus 1fill}\mbox{1829 Sept. . .\quad — 29}\par\smallskip
\noindent\hangindent=1em Die Intensität der magnetischen Kraft der Erde, auf absolutes Maß zurückgeführt . . .\unskip\hspace{1em plus 1fill}\mbox{1832 Dez. . .\quad — 79}\par\smallskip
\noindent\hangindent=1em Allgemeine Theorie des Erdmagnetismus . . . .\unskip\hspace{1em plus 1fill}\mbox{1838\quad 119}\par\smallskip
\noindent\hangindent=1em Allgemeine Lehrsätze über Anziehungs- und Abstoßungskräfte, die umgekehrt proportional zum Quadrat der Entfernung wirken \unskip\hspace{1em plus 1fill}\mbox{1839\quad — 195}\par\smallskip
\noindent\hangindent=1em Dioptrische Untersuchungen . . . . . . . . . . . . . .\unskip\hspace{1em plus 1fill}\mbox{1840 Dez. . .\quad — 243}\par\smallskip
\par\Needspace{5\baselineskip}\medskip\noindent\textit{Anzeigen eigener Abhandlungen.}\par\smallskip
\noindent\hangindent=1em Theorie der Anziehung homogener ellipsoidischer Körper .\unskip\hspace{1em plus 1fill}\mbox{1813 April . .\quad — 279}\par\smallskip
\noindent\hangindent=1em Allgemeine Grundlagen der Theorie der Gestalt von Flüssigkeiten im Gleichgewicht . . . .\unskip\hspace{1em plus 1fill}\mbox{1829 Okt. . .\quad — 287}\par\smallskip
\noindent\hangindent=1em Die Intensität der magnetischen Kraft der Erde, auf absolutes Maß zurückgeführt . . .\unskip\hspace{1em plus 1fill}\mbox{1832 Dez. . .\quad — 293}\par\smallskip
\noindent\hangindent=1em Allgemeine Lehrsätze in Bezug auf das umgekehrte Verhältnis usw. \unskip\hspace{1em plus 1fill}\mbox{1840 März .\quad — 305}\par\smallskip
\noindent\hangindent=1em Dioptrische Untersuchungen . . . . . . . . . . . . . . .\unskip\hspace{1em plus 1fill}\mbox{1841 Jan. . .\quad — 309}\par\smallskip
\par\Needspace{5\baselineskip}\medskip\noindent\textit{Verschiedene Aufsätze über Magnetismus.}\par\smallskip
\noindent\hangindent=1em Erdmagnetismus und Magnetometer . . . . . . . . . .\unskip\hspace{1em plus 1fill}\mbox{1836\quad — 315}\par\smallskip
\noindent\hangindent=1em Einleitung für die Zeitschrift: Resultate usw. . . . . . . . .\unskip\hspace{1em plus 1fill}\mbox{1836\quad — 345}\par\smallskip
\noindent\hangindent=1em Ein neues Hilfsmittel für die magnetischen Beobachtungen . . . . . .\unskip\hspace{1em plus 1fill}\mbox{1837 Okt\quad — 352}\par\smallskip
\noindent\hangindent=1em Über ein neues Instrument, das zunächst zur unmittelbaren Beobachtung von Änderungen der Intensität der horizontalen Komponente des Erdmagnetismus bestimmt ist . . . . . . . . . . . . . . . . .\unskip\hspace{1em plus 1fill}\mbox{1837\quad — 357}\par\smallskip
\noindent\hangindent=1em Anleitung zur Bestimmung der Schwingungsdauer einer Magnetnadel . .\unskip\hspace{1em plus 1fill}\mbox{1837\quad — 374}\par\smallskip
\noindent\hangindent=1em Über ein Mittel, die Beobachtung von Ablenkungen zu erleichtern . .\unskip\hspace{1em plus 1fill}\mbox{1839\quad — 395}\par\smallskip
\noindent\hangindent=1em Zur Bestimmung der Konstanten des Bifilarmagnetometers . . . . .\unskip\hspace{1em plus 1fill}\mbox{1840\quad — 404}\par\smallskip
\noindent\hangindent=1em Vorschriften zur Bestimmung der magnetischen Wirkung, die ein Magnetstab in der Ferne ausübt . . . . . . . . . . . . . .\unskip\hspace{1em plus 1fill}\mbox{1840\quad — 427}\par\smallskip
\SourceLeaf{648}
\noindent\hangindent=1em Über die Anwendung des Magnetometers zur Bestimmung der absoluten Deklination . . . . . . . . . . . . . . . . . . . .\unskip\hspace{1em plus 1fill}\mbox{1841\quad Seite 436}\par\smallskip
\noindent\hangindent=1em Beobachtungen der magnetischen Inklination in Göttingen . . . . . .\unskip\hspace{1em plus 1fill}\mbox{1841\quad — 444}\par\smallskip
\par\Needspace{5\baselineskip}\medskip\noindent\textit{Aufsätze über verschiedene Themen der mathematischen Physik.}\par\smallskip
\noindent\hangindent=1em Fundamentalgleichungen für die Bewegung schwerer Körper auf der rotierenden Erde . . . . . . . . . . . . . . . . . . . .\unskip\hspace{1em plus 1fill}\mbox{1803\quad — 495}\par\smallskip
\noindent\hangindent=1em Über achromatische Doppelobjektive, insbesondere im Hinblick auf die vollständigere Aufhebung der Farbzerstreuung . . . . . . . .\unskip\hspace{1em plus 1fill}\mbox{1817 Dez. . .\quad — 504}\par\smallskip
\noindent\hangindent=1em Brief an \textsc{Brandes} über denselben Gegenstand . . . . . . . . . .\unskip\hspace{1em plus 1fill}\mbox{1831\quad — 509}\par\smallskip
\noindent\hangindent=1em Berichtigung der Stellung der Schneiden einer Waage . . . . . . .\unskip\hspace{1em plus 1fill}\mbox{1837 März . .\quad — 511}\par\smallskip
\par\Needspace{5\baselineskip}\medskip\noindent\textit{Physikalische Beobachtungen.}\par\smallskip
\noindent\hangindent=1em Nordlicht am 7. Januar 1831 . . . . . . . . . . . . . . .\unskip\hspace{1em plus 1fill}\mbox{1831 Febr. . .\quad — 517}\par\smallskip
\noindent\hangindent=1em Magnetisches Observatorium in Göttingen . . . . . . . . . . .\unskip\hspace{1em plus 1fill}\mbox{1834 Aug. . .\quad — 519}\par\smallskip
\noindent\hangindent=1em Beobachtungen der magnetischen Variation in Göttingen und Leipzig . .\unskip\hspace{1em plus 1fill}\mbox{1834 Okt. . .\quad — 523}\par\smallskip
\noindent\hangindent=1em Magnetisches Observatorium in Göttingen . . . . . . . . . . .\unskip\hspace{1em plus 1fill}\mbox{1835 März . .\quad — 528}\par\smallskip
\noindent\hangindent=1em Beobachtungen der magnetischen Variation in Kopenhagen und Mailand .\unskip\hspace{1em plus 1fill}\mbox{1835 März . .\quad — 537}\par\smallskip
\noindent\hangindent=1em Magnetisches Observatorium in Göttingen . . . . . . . . . . .\unskip\hspace{1em plus 1fill}\mbox{1836 Juni . .\quad — 540}\par\smallskip
\noindent\hangindent=1em Das bei den Beobachtungsterminen anzuwendende Verfahren . . . . .\unskip\hspace{1em plus 1fill}\mbox{1836\quad — 541}\par\smallskip
\noindent\hangindent=1em Auszug aus dreijährigen täglichen Beobachtungen der magnetischen Deklination in Göttingen . . . . . . . . . . . . . . . . .\unskip\hspace{1em plus 1fill}\mbox{1836\quad — 556}\par\smallskip
\noindent\hangindent=1em Erläuterungen zu den Terminszeichnungen und den Beobachtungszahlen .\unskip\hspace{1em plus 1fill}\mbox{1836\quad — 568}\par\smallskip
\noindent\hangindent=1em Erläuterungen zu den Terminszeichnungen und den Beobachtungszahlen .\unskip\hspace{1em plus 1fill}\mbox{1837\quad — 576}\par\smallskip
\noindent\hangindent=1em Der magnetische Südpol der Erde . . . . . . . . . . . . .\unskip\hspace{1em plus 1fill}\mbox{1841\quad — 580}\par\smallskip
\par\Needspace{5\baselineskip}\medskip\noindent\textit{Anzeigen fremder Schriften.}\par\smallskip
\noindent\hangindent=1em \textsc{Benzenberg.} Über die \textsc{Dalton}sche Theorie . . . . . . . . . .\unskip\hspace{1em plus 1fill}\mbox{1830 Dez. . .\quad — 583}\par\smallskip
\noindent\hangindent=1em \textsc{Fischer.} Künstliche Magnete . . . . . . . . . . . . . . .\unskip\hspace{1em plus 1fill}\mbox{1832 Sept. .\quad — 591}\par\smallskip
\noindent\hangindent=1em Ergebnisse der Beobachtungen des Magnetischen Vereins . . . . .\unskip\hspace{1em plus 1fill}\mbox{1837 Juni . .\quad — 595}\par\smallskip
\noindent\hangindent=1em \textsc{Gerling.} Einrichtung des mathematisch-physikalischen Instituts . . .\unskip\hspace{1em plus 1fill}\mbox{1842 April .\quad — 596}\par\smallskip
\par\Needspace{9\baselineskip}\medskip\noindent\textit{Nachlass.}\par\smallskip
\noindent\hangindent=1em Zur Elektrodynamik . . . . . . . . . . . . . . . . . . . . . . . .\unskip\hspace{1em plus 1fill}\mbox{\quad — 601}\par\smallskip
\noindent\hangindent=1em Über Kugelfunktionen . . . . . . . . . . . . . . . . . . . . . . . .\unskip\hspace{1em plus 1fill}\mbox{\quad — 630}\par\smallskip
\noindent\hangindent=1em Zum Gebrauch des Komparators . . . . . . . . . . . . . . . . . . . .\unskip\hspace{1em plus 1fill}\mbox{\quad — 632}\par\smallskip
\noindent\hangindent=1em Allgemeine Formeln für die Wirkung eines leuchtenden Punkts $P$ auf einen Punkt $p$ . .\unskip\hspace{1em plus 1fill}\mbox{\quad — 635}\par\smallskip
\noindent\hangindent=1em Bemerkungen . . . . . . . . . . . . . . . . . . . . . . . . . . . .\unskip\hspace{1em plus 1fill}\mbox{\quad — 637}\par\smallskip
\par\medskip\noindent Lithografische Tafel zur Theorie des Erdmagnetismus, Seite 176.
\endgroup

\par\nopagebreak[4]\medskip\noindent\makebox[\linewidth][c]{\rule{.15\linewidth}{.35pt}}\par\nopagebreak[4]
\bigskip
\begin{center}
GÖTTINGEN,\par
\small GEDRUCKT IN DER DIETERICHSCHEN UNIVERSITÄTS-DRUCKEREI\par
W. FR. KAESTNER.
\end{center}
